\subsection{A complete large-conductor packet and all ordinary genus-two classes over \texorpdfstring{$\F_4$}{F4}}

\begin{claim}{A.C2.25}{Proved}
Let $\alpha^2+\alpha+1=0$. The smooth proper curve
\[
 C_0/\F_4:\qquad
 v^2+v=z+\frac{\alpha^2(z+1)}{z^2+z+\alpha}                 \tag{C2.78}
\]
and its coefficientwise-square conjugate satisfy A1.
Consequently all $48$ geometric ordinary genus-two
moduli classes fixed by fourth-power Frobenius
satisfy A1. In particular this holds for every
ordinary genus-two curve over $\F_4$, over every
finite extension, for every finitely generated
form of any of these geometric classes, and
for every hyperbolic open subcurve.

The additional two classes are certified by a
connected cyclic \'etale cover of odd degree
\[
 M=7\cdot31\cdot2503\cdot4519=2\,454\,499\,369             \tag{C2.79}
\]
whose primitive conductor-$M$ Jacobian factor
has two-rank zero. Universal A1 remains A.30.
\end{claim}
\begin{proof}
\proofparagraph{The curve and the exact theta origin.}
Set
\[
 h=z^2+z+\alpha,\qquad G=z^3+z^2+z+\alpha^2,
 \qquad Y=hv,
\]
so that the polynomial model is
\[
                         Y^2+hY=f,\qquad f=hG.             \tag{C2.80}
\]
Here $f$ is monic of degree five. At a root
of $h$, the value $G=\alpha^2(z+1)$ is nonzero;
thus the affine model is smooth. The original
Artin--Schreier equation has three simple poles,
the two roots of $h$ and infinity, and defines
an ordinary genus-two curve. Write $W_1,W_2$
for its finite ramification points and $W_\infty$
for its point at infinity. The different and
the divisor of $dz$ give
\[
 \operatorname{div}(dz)=2(W_1+W_2-W_\infty),\qquad
 B_\Theta=\mathcal O(W_1+W_2-W_\infty).                     \tag{C2.81}
\]
The locally exact differential $dz$ identifies
this with the Raynaud bundle in A.C2.5. Hence
\[
 \Theta=\{[R-B_\Theta]:R\in C_0(\overline{\F}_2)\}.
\]
The pair $T=(h,0)$ represents
$[W_1+W_2-2W_\infty]$. Its double is zero,
since $\operatorname{div}(h)=2W_1+2W_2-4W_\infty$.
Consequently
\[
                  [R-W_\infty]+T=[R-B_\Theta].             \tag{C2.82}
\]
This fixes both the origin and the sign of the
theta test. The individual finite branch points
need not be rational over the certificate field;
their sum and $T$ are rational.

Choose $r$ with $r^2+r=\alpha$ and put $x=z+r$.
After removing a constant, (C2.78) has the
normalized coefficient triple
\[
                 (1,\alpha^2(r+1),\alpha^2r).
\]
Its symmetric functions are $\alpha,0,\alpha^2$,
so it is key $203$ in A.C2.22. Squaring the
coefficients gives key $302$. This field
automorphism transports the theta classes,
all algebraic identities below, and the
associated cyclic cover. It therefore suffices
to prove the assertion for $C_0$.

\proofparagraph{The fixed field and exact torsion generator.}
All computations take place in
\[
 F=\F_2[w]/(w^{18}+w^7+1)=\F_{4^9}.                         \tag{C2.83}
\]
An integer between $0$ and $2^{18}-1$ encodes
the polynomial whose binary digits are its
coefficients in ascending powers of $w$.
The binary modulus is $262273$, the primitive
generator is $2$, and $\alpha$ is $114587$.
The verifier checks the complete nonzero
power cycle of length $262143$. An independent
verifier also checks the polynomial
irreducibility criterion at exponents $18,9,6$.

The immutable data file
\path{code/A_F4_large_theta_witness.json}
contains a theta generator
\[
 R=(49,185454),\qquad \eta=[R-B_\Theta],
\]
whose reduced Mumford pair on (C2.80) is
\[
 U=z^2+124780z+171897,
 \qquad V=26069z+117672.                                   \tag{C2.84}
\]
The coordinates of $R$ are the original $(z,v)$
coordinates, whereas $V$ is the $Y$ coordinate
of the reduced divisor. Every displayed
coefficient is interpreted in (C2.83).
The exact checks are
\[
 [M]\eta=0,\qquad [M/p]\eta\ne0
       \quad(p=7,31,2503,4519).                             \tag{C2.85}
\]
The four factors are verified prime by trial
division, so $\eta$ has exact order $M$.

For an additional check of the primary module,
put
\[
 R_H=(1,114586),\qquad H=[4M/7]\,[R_H-B_\Theta].
\]
Its reduced pair is
\[
                   (z^2+z+1,\,114586z).
\]
The verifier proves $[7]H=0$ and
$H\ne[jM/7]\eta$ for $0\le j\le6$.
Thus it explicitly proves two independent
seven-torsion classes.
The actual $\F_4$ Frobenius polynomial from
A.C2.24 gives
\[
 \begin{aligned}
 P_4(X)&=X^4+2X^3+X^2+8X+16,\\
 \#J(F)&=68\,725\,982\,332
      =4\cdot7^2\cdot31\cdot2503\cdot4519,\\
 \#C_0(F)&=262171.
 \end{aligned}                                             \tag{C2.86}
\]
One verifier uses $\det(\pi^9-I)$ on the
integral companion matrix; two others use
integer Newton recurrences. These are the
curve and Jacobian Frobenius formulas of
\citet[II, Theorem 1.1; III, Theorem 11.1
and Corollary 11.4]{A_Milne2008}.
The group order and the independent
seven-torsion classes prove
\[
 J(F)_{\rm odd}\cong
 C_7^2\oplus C_{31}\oplus C_{2503}\oplus C_{4519}.
\]
Thus $M$ is its exact exponent. The main
certificate needs only (C2.85) and the
point-multiple identities below; this
additional argument justifies the module
structure used in the discovery search.

\proofparagraph{Divisor meaning of the finite arithmetic.}
A reduced pair $(U,V)$ means that $U$ is monic,
$\deg V<\deg U\le2$, and $U\mid V^2+hV+f$.
It represents $D(U,V)-\deg(U)W_\infty$,
where $D(U,V)$ is the finite divisor of the
ideal $(U,Y-V)$. This representation is
unique; zero is $(1,0)$
\citep[Theorem 2.2, author p.~5]{A_Lange2005}.
The checker implements the general
composition and reduction rules, including
coincident points and branch fibres.

For two pairs, compute the monic extended gcds
\[
 \begin{aligned}
 d_1&=a_0U_1+b_0U_2=\gcd(U_1,U_2),\\
 d&=e d_1+c(V_1+V_2+h)
       =\gcd(d_1,V_1+V_2+h),\qquad a=ea_0,\quad b=eb_0.
 \end{aligned}
\]
Composition is
\[
 U=\frac{U_1U_2}{d^2},\qquad
 V=\frac{aU_1V_2+bU_2V_1+c(V_1V_2+f)}d\pmod U.             \tag{C2.87}
\]
While $\deg U>2$, reduction replaces it by
\[
 U'=\frac{V^2+hV+f}U,
 \qquad V'=V+h\pmod {U'},                                  \tag{C2.88}
\]
making each new first polynomial monic.
These are the characteristic-two forms of
\citet[Algorithms 3.1--3.2, author pp.~5--6]
{A_Lange2005}, which explicitly allow the
general $h$ polynomial and even characteristic.

To recall why these are divisor identities,
the product of the two ideals is generated by
\[
 \begin{gathered}
 U_1U_2,\quad U_1(Y+V_2),\quad U_2(Y+V_1),\\
 (V_1+V_2+h)Y+(V_1V_2+f).
 \end{gathered}
\]
The Bezout relation produces $d(Y+V)$.
Locally, $d$ removes exactly the common
opposite-point pairs, giving
$I_1I_2=(d)I(U,V)$. A repeated ramification
point is removed using
$\operatorname{div}(z-z_0)=2W-2W_\infty$.
Thus (C2.87) represents the sum of classes.
For (C2.88), the norm of $Y+V$ is
$V^2+hV+f=UU'$. Its finite zero divisor is
$D(U,V)+D(U',V)$, with multiplicities, and
its only pole is at infinity. The two
degree-zero divisor classes are negatives.
Hyperelliptic conjugation changes $V$ to
$V+h$, proving the reduction rule. Monic
normalization changes only a nonzero scalar.

Every division is checked to have zero
remainder, and the defining congruence is
checked after every operation. Binary
repeated addition therefore supplies exact
identities in $\operatorname{Pic}^0(C_0)$,
not approximate numerical relations.
For a saved point $(x,y)$, the checker starts
with $([x,1],[h(x)y])$, then adds $(h,0)$.
By (C2.82) the resulting class is
$[R-B_\Theta]$.

\proofparagraph{Complete primitive orbit coverage.}
The fixed data contain $1296$ records
$(x_i,y_i,i)$ such that
\[
 \gcd(i,M)=1,
 \qquad [i]\eta=[R_i-B_\Theta]\in\Theta.                    \tag{C2.89}
\]
Every record, not a sample, is verified by
the preceding exact arithmetic. The local
orders of two at $7,31,2503,4519$ are
$3,5,1251,753$. Hence
\[
 \operatorname{ord}_M(2)=1\,570\,005,
 \qquad \varphi(M)=2\,034\,726\,480,
 \qquad \frac{\varphi(M)}{\operatorname{ord}_M(2)}=1296.     \tag{C2.90}
\]
It remains to prove that the recorded indices
belong to distinct doubling orbits.

One verification uses the primitive root
$3$ at each prime. Integer discrete logarithms
identify the unit group with
$\Z/6\oplus\Z/30\oplus\Z/2502\oplus\Z/4518$.
The lattice generated by these four period
vectors and the log vector of two has column
Hermite basis
\[
 \begin{pmatrix}
 6&0&0&2\\0&6&0&0\\0&0&6&4\\0&0&0&6
 \end{pmatrix},\qquad \det=1296.                            \tag{C2.91}
\]
Reducing from the last column to the first
gives $1296$ different residues for the
saved indices. The vector of two reduces
to zero. This proves exhaustion.

A separate verification avoids discrete-log
normalization and Hermite normal form.
Enumerate the local doubling cycles, and
label each local unit by its smallest cycle
representative and its time coordinate $t_p$.
Two global units are in one orbit exactly
when their local representatives agree and
the required time differences admit one
common shift. By the generalized Chinese
remainder theorem this is equivalent to
agreement of every pair of shifts modulo
the gcd of the two local periods. The four
local representatives and all such relative
times give $1296$ distinct signatures on
the data. A third implementation uses
different primitive roots and the equivalent
signature with local log residues modulo
$(2,6,2,6)$ and two relative times modulo
three; its $144\cdot9$ signatures also exhaust
the quotient. These are independent exact
checks of the inference from (C2.89) to
all primitive doubling orbits.

\proofparagraph{The cover and the section conclusion.}
Let $L$ represent $\eta$ over an algebraic
closure and choose $L^M\simeq\mathcal O$.
The algebra $\bigoplus_{j=0}^{M-1}L^{-j}$
defines a finite \'etale $\mu_M$-torsor
$D\to C_0$, since $M$ is odd. It is
geometrically connected because $L$ has
exact order $M$. After a finite constant
extension its deck group is cyclic of
order $M$, and its genus is $M+1$.

For every nonzero character $i$,
$H^1(C_0,L^i)$ is one dimensional. The
exact sequence
\[
 0\longrightarrow L^i\longrightarrow F_*L^{2i}
 \longrightarrow B_\Theta\otimes L^i\longrightarrow0
\]
makes its Frobenius arrow zero exactly
when $L^i\in\Theta$, as in A.C2.5.
Every primitive cycle contains a zero arrow
by (C2.89)--(C2.91). Frobenius is therefore
nilpotent on the entire primitive
conductor-$M$ summand, and that Jacobian
factor has two-rank zero. Lower-conductor
characters need not vanish.

A.C2.19, or equivalently the all-genus
criterion A.C2.20, now proves A1 for $C_0$.
Four distinct prime factors are permitted:
the primitive operator has support bound
$2^4=16$, and $M>15^2$, exactly as required
in that proof. This invokes the proved
geometric divisor argument there, not an
inference from the size of the cover.

The property is geometric. For a form over
a finitely generated field, the isomorphism,
the torsion class, its trivialization and
the deck maps are defined after a finite
extension of the whole base field. The
geometric formulation in A.C2.19--A.C2.20,
together with A.3 and the descent in A.8,
applies. An original collision would persist
over that extension. Functoriality to the
compactification handles open subcurves.
Coefficient conjugation gives key $302$,
and A.C2.22 supplies the other $46$ classes.

\proofparagraph{Reproducible certificates and exact scope.}
The fixed witness file has SHA256
\begin{center}
\path{edcf6ed3d425eb72837ae72f521f7221ee91c8a53c1b005885a4c1280d530ae6}.
\end{center}
The canonical checker
\path{code/A_F4_large_theta_canonical.py}
and its adjacent JSON verify all $1296$
relations and the order statements in
$60773$ exact divisor operations. It
imports no C implementation, search code,
accelerated polynomial engine, or tailored
coset normalization. Root and the independent
reviewer each reran it successfully.
The further independent programs
\path{code/A_F4_large_theta_root_orbits.py}
and
\path{code/A_F4_large_theta_degeneration.py},
with their JSON outputs, verify the field,
point equations, group order, reindexing
and complete orbit signatures by the
distinct methods described above.
The full mathematical reviews are
\path{audit/A_F4_large_theta_degeneration_review.md}
and
\path{audit/A_F4_all_48_root_review.md}.

The discovery code and full data are
retained under
\path{code/A_F4_large_theta_search.py/.json}.
Discovery stopped when this positive
packet was complete; no negative
enumeration is needed. A.C2.24's exact
small-primary and mixed-support exclusions
remain valid, including the entire
$73$-primary obstruction. They do not
exclude the different mixed conductor
(C2.79). The fixed-prime obstruction
A.C2.14 and the abelian-tower descent
A.C2.23 are also unchanged. This proves
the stated finite moduli class theorem;
universal A1 remains unresolved.
\end{proof}
